\subsection{Some operations in the tensor algebra.}

In this section we shall denote by \(e_x\) ( \(x \in E\) ) the linear mapping \(u \to x \otimes u\) of the tensor algebra T(E) into itself.

Lemma 1. — Let f be an element of the dual E* of E. There exists one and only one linear mapping \(i_{f}\) of T(E) into itself such that

\begin{enumerate}
\def\labelenumi{(\arabic{enumi})}
\setcounter{enumi}{3}
\tightlist
\item
\end{enumerate}

\[
i _ {f} (1) = 0\tag{5}
\]

\[
i _ {f} \circ e _ {x} + e _ {x} \circ i _ {f} = f (x). I
\]

for all \(x \in \mathbf{E}\)

(where I denotes the identity mapping). The mapping \(f \to i_f\) of E* into \(\mathfrak{U}(\mathrm{T}(\mathrm{E}))\) is linear. One has \(i_f(\mathrm{T}^n) \subset \mathrm{T}^{n-1}\) , \((i_f)^2 = 0\) , and \(i_f \circ i_g + i_g \circ i_f = 0\) for \(f\) , \(g\) in E*. The mapping \(i_f\) is zero on the subalgebra of

T(E) generated by the kernel of f. The ideal I(Q) is stable under \(i_{f}\) ; by passage to the quotient \(i_{f}\) therefore defines a linear mapping (still denoted \(i_{f}\) ) of C(Q) into itself.

Indeed, formula (5) can be written

\[
i _ {f} (x \otimes u) = - x \otimes i _ {f} (u) + f (x). u \quad (x \in \mathrm{E}, u \in \mathrm{T(E)}).\tag{6}
\]

Since (4) completely determines \(i_{f}\) on \(\mathbf{T}^{0}\) , and since (6) determines \(i_{f}\) on \(\mathbf{T}^{n}\) once its values on \(\mathbf{T}^{n-1}\) are known, the uniqueness of \(i_{f}\) is proved. On the other hand, for \(x\in\mathbf{E}\) and \(u\in\mathbf{T}^{n-1}\) , the right-hand side of (6) is bilinear on \(\mathbf{E}\times\mathbf{T}^{n-1}\) ; this proves the existence of \(i_{f}\) by induction on \(n\) (chap.~III, § 1, no. 2) and also proves, by induction on \(n\) , that \(i_{f}(\mathbf{T}^{n})\subset\mathbf{T}^{n-1}\) . If \(f=ag+bh\) ( \(a,b\) in A, \(g,h\) in E*) it is clear that \(ai_{g}+bi_{h}\) satisfies (4) and (5), hence is equal to \(i_{f}\) . We have

\[
(i _ {f}) ^ {2} \circ e _ {x} = - i _ {f} \circ e _ {x} \circ i _ {f} + f (x) i _ {f} = e _ {x} \circ (i _ {f}) ^ {2}
\]

and, since \((i_f)^2(1) = 0\) we deduce by induction on \(n\) that \((i_f)^2\) is zero on \(\mathbf{T}^n\) . Replacing \(f\) by \(f + g\) , the relation \((i_f)^2 = 0\) gives \(i_f \circ i_g + i_g \circ i_f = 0\) . An induction argument analogous to the preceding ones shows that \(i_f\) is zero on the subalgebra generated by the kernel of \(f\) . Finally, (6) implies that the set of elements \(u\) of I(Q) such that \(i_f(u) \in \mathrm{I}(Q)\) is a left ideal of T(E); moreover, if \(u = (x \otimes x - Q(x).1) \otimes v (x \in E, v \in T(E))\) , one has

\[
\begin{array}{r l} i _ {f} (u) & = f (x) x \otimes v - x \otimes i _ {f} (x \otimes v) - \mathrm{Q} (x) i _ {f} (v) \\ & = f (x) x \otimes v - f (x) x \otimes v + x \otimes x \otimes i _ {f} (v) - \mathrm{Q} (x) i _ {f} (v) \\ & = (x \otimes x - \mathrm{Q} (x). 1) \otimes i _ {f} (v); \end{array}
\]

therefore I(Q) is stable under \(i_{f}\) and the last assertion follows. QED.

Let F be a bilinear form on E; in the rest of this paragraph we shall denote by \(i_x^{\mathrm{F}}\) ( \(x \in \mathrm{E}\) ) the map \(i_f\) corresponding to the linear form \(f: y \to \mathrm{F}(x, y)\) on E.

Lemma 2. --- There exists a linear map \(\lambda_{\mathrm{F}}\) of T(E) into itself such that

\begin{enumerate}
\def\labelenumi{(\arabic{enumi})}
\setcounter{enumi}{6}
\tightlist
\item
\end{enumerate}

\[
\lambda_ {\mathrm{F}} (1) = 1\tag{8}
\]

\[
\lambda_ {\mathrm{F}} \circ e _ {x} = (e _ {x} + i _ {x} ^ {\mathrm{F}}) \circ \lambda_ {\mathrm{F}}\tag{\( (x \in E) \) .}
\]

Whatever the \(f \in \mathbf{E}^*\) , one has

\[
\lambda_ {\mathrm{F}} \circ i _ {f} = i _ {f} \circ \lambda_ {\mathrm{F}}.\tag{9}
\]

Indeed, formula (8) is equivalent to

\[
\lambda_ {\mathbb {F}} (x \otimes u) = x \otimes \lambda_ {\mathbb {F}} (u) + i _ {x} ^ {\mathrm{F}} (\lambda_ {\mathbb {F}} (u)) \quad (x \in \mathrm{E}, u \in \mathrm{T} (\mathrm{E})).\tag{10}
\]

Since (7) determines entirely \(\lambda_{\mathbb{F}}\) on \(\mathrm{T}^{0}\) , and since (10) determines \(\lambda_{\mathbb{F}}\) on \(\mathrm{T}^{n}\) once its values on \(\mathrm{T}^{n-1}\) are known, the uniqueness of \(\lambda_{\mathbb{F}}\) is proved. On the other hand, for \(x \in \mathrm{E}\) and \(u \in \mathrm{T}^{n-1}\) the right-hand side of (10) is bilinear on \(\mathrm{E} \times \mathrm{T}^{n-1}\) ; this proves the existence of \(\lambda_{\mathbb{F}}\) by induction on \(n\) . It remains to prove (9), which we shall do by induction; both sides of (9) vanish on \(\mathrm{T}^{0}\) ; suppose (9) holds on \(\sum_{h=0}^{n-1} \mathrm{T}^{h}\) ; we then have, for \(x \in \mathrm{E}\) and \(u \in \mathrm{T}^{n-1}\) :

\[
\begin{array}{r l} (\lambda_ {\mathbb {F}} \circ i _ {f}) (x \otimes u) & = (- \lambda_ {\mathbb {F}} \circ e _ {x} \circ i _ {f} + f (x) \lambda_ {\mathbb {F}}) (u) \\ & = - (e _ {x} + i _ {x} ^ {\mathrm{F}}) \circ \lambda_ {\mathbb {F}} \circ i _ {f} (u) + f (x) \lambda_ {\mathbb {F}} (u) \\ & = - (e _ {x} + i _ {x} ^ {\mathrm{F}}) \circ i _ {f} \circ \lambda_ {\mathbb {F}} (u) + f (x) \lambda_ {\mathbb {F}} (u) \\ & = (i _ {f} \circ e _ {x} \circ \lambda_ {\mathbb {F}} - f (x) \lambda_ {\mathbb {F}} + i _ {f} \circ i _ {x} ^ {\mathrm{F}} \circ \lambda_ {\mathbb {F}} + f (x) \lambda_ {\mathbb {F}}) (u) \\ & = (i _ {f} \circ (e _ {x} + i _ {x} ^ {\mathrm{F}}) \circ \lambda_ {\mathbb {F}}) (u) = (i _ {f} \circ \lambda_ {\mathbb {F}}) (x \otimes u), \end{array}
\]

whence our last assertion.

Lemma 3. --- Let F and G be two bilinear forms on E. Then we have \(\lambda_{\mathrm{F}} \circ \lambda_{\mathrm{G}} = \lambda_{\mathrm{F+G}}\) . For every bilinear form F on E, \(\lambda_{\mathrm{F}}\) is a bijection from T(E) onto itself.

Indeed, \(\lambda_{\mathbb{F}} \circ \lambda_{\mathbb{G}}\) possesses the characteristic properties (7) and (8) of \(\lambda_{\mathbb{F} + \mathbb{G}}: \text{one has } (\lambda_{\mathbb{F}} \circ \lambda_{\mathbb{G}})(1) = 1\) , and

\[
\begin{array}{r l} \lambda_ {\mathrm{F}} \circ \lambda_ {\mathrm{G}} \circ e _ {x} & = \lambda_ {\mathrm{F}} \circ (e _ {x} + i _ {x} ^ {\mathrm{G}}) \circ \lambda_ {\mathrm{G}} = (e _ {x} + i _ {x} ^ {\mathrm{G}} + i _ {x} ^ {\mathrm{F}}) \circ \lambda_ {\mathrm{F}} \circ \lambda_ {\mathrm{G}} \\ & = (e _ {x} + i _ {x} ^ {\mathrm{F+G}}) \circ \lambda_ {\mathrm{F}} \circ \lambda_ {\mathrm{G}}. \end{array}
\]

On the other hand, if F = 0, we have \(i_x^{\mathbb{F}} = 0\) for all \(x \in \mathbb{E}\) and consequently \(\lambda_{\mathbb{F}} = \mathbb{I}\) which implies that, for every F, we have \(\lambda_{\mathbb{P}} \circ \lambda_{-\mathbb{P}} = \lambda_{-\mathbb{F}} \circ \lambda_{\mathbb{F}} = \mathbb{I}\) .

